\documentclass[11pt]{article}
\usepackage[margin=1in]{geometry}
\usepackage{amsmath}
\usepackage{booktabs}
\usepackage{hyperref}
\usepackage{longtable}
\usepackage{array}

\title{War Carbon Emissions Dashboard\\
  \large Methodology v1.0.5\\[0.4em]
  \normalsize\normalfont Working Draft --- Subject to Revision}
\author{WCED Project}
\date{Last updated: 2026-05-24}

\begin{document}
\maketitle

\noindent\emph{This document specifies the current emission estimation
methodology used by the War Carbon Emissions Dashboard. Equations, parameters,
and decision rules are subject to revision; each revision increments the
methodology version and is documented in the accompanying changelog
(Section~\ref{sec:versions}). An academic preprint with motivation, related
work, and results is in preparation.}

\vspace{1em}

%-----------------------------------------------------------------------
\section{Scope}
%-----------------------------------------------------------------------

WCED v1.0.5 quantifies $\mathrm{CO_2}$ released by hydrocarbon combustion
at oil and fuel infrastructure in the Iran--Gulf theatre of the 2026
Iran--Israel--US war. This document specifies the equations, parameter
priors, verification rules, and reconciliation logic used by the estimation
system.

Version~1 covers refineries, crude and product depots, and petrochemical
complexes from 2026-02-28 onward. It does \emph{not} cover
destroyed-building embodied carbon, munitions embodied carbon, combat
aviation fuel, shipping rerouting, reconstruction, or non-carbon ecological
damage; each of those is deferred to a later version. Emission estimates
represent $\mathrm{CO_2}$ released from observed fires and are not
attributed to specific belligerents.

%-----------------------------------------------------------------------
\section{Estimation Framework}
%-----------------------------------------------------------------------

All emission estimates are expressed as probability distributions
characterised by a 5th~percentile (p5), median (p50), and 95th~percentile
(p95), computed by Monte Carlo simulation. Every result carries a
provenance record tracing it through a chain of inputs back to primary
satellite observations; no estimate is published as a scalar without bounds.

Each fire event is assigned a confidence label that reflects the evidence
supporting it. Labels are defined in Section~\ref{sec:verify} and propagate
forward into every downstream emission estimate, so the uncertainty
associated with event identification is never decoupled from the reported
emission range.

The two core emission methods---Fire Radiative Power (FRP)-based
(Section~\ref{sec:frp}) and inventory-based
(Section~\ref{sec:inventory})---are computed independently and then
reconciled (Section~\ref{sec:reconcile}). All emission factors and
parameter prior distributions are maintained in versioned reference tables
external to the estimation logic; they are not hard-coded.

%-----------------------------------------------------------------------
\section{Emission Calculations}
%-----------------------------------------------------------------------

\subsection*{Sections 3.1--3.2\enspace(reserved)}
Facility registry schema and emission factor tables are maintained as
versioned reference data and are not reproduced here.

%-----------------------------------------------------------------------
\subsection{FRP-based Method}
\label{sec:frp}
%-----------------------------------------------------------------------

Fire Radiative Power (FRP) observations from VIIRS, MODIS, and Sentinel-3
SLSTR are clustered into an event-level time series from which a total fire
radiative energy is derived. The method proceeds in four stages: raw energy
integration, baseline subtraction, overpass-gap correction, and conversion
to $\mathrm{CO_2}$ mass.

\paragraph{Stage 1: Raw energy integration.}
The pipeline begins from satellite FRP observations sampled at overpass
times. The raw Fire Radiative Energy is the trapezoidal integral over the
observation window:
\begin{equation}
I_{\mathrm{raw}} \;=\; \int_{t_0}^{t_1} \mathrm{FRP}(t)\,\mathrm{d}t.
\label{eq:fre-raw}
\end{equation}

\paragraph{Stage 2: Pre-war baseline subtraction.}
The raw integral over-counts because operating facilities flare routinely.
A facility-specific pre-war baseline is subtracted to isolate the
conflict-attributable signal. The baseline is derived from 12 months of
archival satellite observations preceding the conflict
(2025-02-28 to 2026-02-27).

For each facility, the rolling 30-day pre-war baseline FRP and its
uncertainty are:
\begin{align}
\mathrm{FRP}_{\mathrm{baseline}} &=
  P_{75}\!\bigl(\{\mathrm{FRP}_i\}_{i \in W_{30}}\bigr),
\label{eq:baseline} \\
\sigma_{\mathrm{baseline}} &=
  \mathrm{IQR}\!\bigl(\{\mathrm{FRP}_i\}\bigr) \big/ 1.349,
\label{eq:baseline-std}
\end{align}
where $W_{30}$ is the 30-day window preceding the event, $P_{75}$ denotes
the 75th percentile (capturing the upper envelope of routine operations),
and the IQR-based estimator provides a robust standard deviation resistant
to transient flare-ups. The net energy available for conflict attribution
is:
\begin{equation}
I_{\mathrm{net}} = \max\!\bigl(0,\;
  I_{\mathrm{raw}} - \mathrm{FRP}_{\mathrm{baseline}} \cdot \Delta t\bigr),
\label{eq:net-frp}
\end{equation}
where $\Delta t$ is the event duration in days. Events for which the
30-day window contains fewer than five observations use a conservative
fallback baseline ($\mu = 0$~MW, $\sigma = 50$~MW) and are flagged as
having insufficient baseline history.

\paragraph{Stage 3: Overpass-gap correction.}
Satellites sample at discrete overpass times. Two multiplicative corrections
account for fire activity in unobserved intervals and for residual sampling
bias. The corrected Fire Radiative Energy is:
\begin{equation}
I_{\mathrm{FRE}} \;=\; k_{\mathrm{ext}} \cdot d \cdot I_{\mathrm{net}},
\label{eq:fre}
\end{equation}
where $d$ is the burn duty cycle, representing the fraction of time the
fire is active between overpasses (prior: triangular(0.40, 0.70, 0.95);
informed by the Kuwait 1991 and Abqaiq 2019 oil-facility fires), and
$k_{\mathrm{ext}}$ is an extrapolation factor correcting residual bias
when overpass sampling is sparse (prior: $\mathcal{N}(1.0,\,0.15)$,
calibrated against MODIS/VIIRS cross-comparisons). Both parameters are
sampled independently in the Monte Carlo (Section~\ref{sec:uq}); applying
either as a point value before sampling would systematically narrow the
resulting uncertainty interval. $I_{\mathrm{net}}$ and $I_{\mathrm{FRE}}$
both carry units of MJ.

\paragraph{Stage 4: Conversion to $\mathrm{CO_2}$ mass.}
Corrected energy converts to fuel mass through a calibrated coefficient,
then to $\mathrm{CO_2}$ mass through combustion stoichiometry:
\begin{align}
m_{\mathrm{fuel}} &= \alpha \cdot I_{\mathrm{FRE}},
\label{eq:fuel}\\
m_{\mathrm{CO_2}} &= m_{\mathrm{fuel}} \cdot \beta \cdot r,
\label{eq:co2-frp}
\end{align}
where $\alpha$ is the FRP-to-combustion-rate coefficient
($\mathcal{N}(0.368,\,0.05)$~kg/MJ; Wooster et al.~\cite{wooster2005}),
$\beta = (44/12)\cdot f_C$ is the stoichiometric
$\mathrm{C}\!\to\!\mathrm{CO_2}$ conversion factor with $f_C \approx 0.86$
the carbon mass fraction of crude oil (EPA AP-42; \cite{ap42}), and $r$ is
the carbon recovery fraction (triangular(0.92, 0.96, 0.98);
Hobbs \& Radke~\cite{hobbs1992}). At central prior values,
$\beta \cdot r \approx 3.15$~kg$\,\mathrm{CO_2}$ per kg crude burned,
consistent with the value used in the CCI Iran brief
\cite{bigger2026,cci2026}.

%-----------------------------------------------------------------------
\subsection{Inventory-based Method}
\label{sec:inventory}
%-----------------------------------------------------------------------

For facilities with a curated nameplate capacity $C$ (barrels), a
bottom-up estimate is computed from:
\begin{equation}
m_{\mathrm{CO_2}} \;=\; C \cdot \phi \cdot \psi \cdot \mathrm{EF},
\label{eq:co2-inv}
\end{equation}
where $\phi$ is the fraction of nameplate capacity present at the time of
the strike (prior: uniform(0.30, 0.90); a judgment-based prior pending
empirical calibration), $\psi$ is the destroyed fraction, and
$\mathrm{EF}$ is the combustion emission factor for the relevant product:

\begin{itemize}
\item Crude depots and refinery crude units:
  $\mathrm{EF} = 0.425$~tCO$_2$/barrel
  (EPA AP-42 \S1.3; consistent with CCI Iran brief \cite{cci2026}).
\item Refined-product depots:
  $\mathrm{EF} = 0.430$~tCO$_2$/barrel
  (IPCC 2006 Guidelines, Vol.~2, Ch.~1, Table~1.3; \cite{ipcc2006}).
\end{itemize}

Each emission factor is designated applicable to specific facility types
only; type-factor mismatches are rejected at the estimation stage. Applying
a crude-oil factor to a gas-processing facility would produce
systematically incorrect results and is therefore prohibited.

\paragraph{Facility-type restriction.}
The inventory method applies only to storage-type facilities
(oil depots, storage tank farms, and refineries) where nameplate capacity
represents stored inventory. It does not apply to production or processing
facilities (gas processing plants, petrochemical complexes) where capacity
represents throughput and on-site inventory is a small fraction of stated
capacity. Events at production facilities use the FRP-based method
exclusively.

\paragraph{Destroyed-fraction prior (v1.0.5 recalibration).}
The parameter $\psi$ is event-specific: when a reviewer provides an
explicit assessment $\hat\psi$ from before/after satellite imagery, the
prior centres on that value using triangular$(\hat\psi - 0.15,\;
\hat\psi,\; \hat\psi + 0.15)$, truncated to $[0,1]$. When no assessment
is available, the v1.0.5 defaults in Table~\ref{tab:psi} apply.

The defaults were recalibrated in v1.0.5 after systematic FRP/inventory
agreement ratios of 2.7--7.4$\times$ were observed under earlier defaults,
indicating that initial priors substantially over-estimated the fraction of
storage capacity destroyed in a typical strike:

\begin{table}[h]
\centering
\caption{Fraction-destroyed default priors by facility type and
  methodology version.}
\label{tab:psi}
\begin{tabular}{@{}llll@{}}
\toprule
\textbf{Facility type} & \textbf{Prior v1.0--v1.0.4}
  & \textbf{Prior v1.0.5} & \textbf{Basis} \\
\midrule
Oil depot             & tri(0.20, 0.40, 0.65)
  & tri(0.05, 0.15, 0.30) & FRP/inv ratio calibration \\
Storage tank farm     & tri(0.20, 0.40, 0.65)
  & tri(0.05, 0.15, 0.30) & Same \\
Refinery              & tri(0.10, 0.20, 0.35)
  & tri(0.05, 0.15, 0.30) & Same \\
\bottomrule
\end{tabular}
\end{table}

The recalibrated mode of 0.15 reflects that strikes typically produce
localised tank fires rather than facility-wide destruction, consistent
with post-strike optical satellite imagery.

%-----------------------------------------------------------------------
\subsection{Reconciliation}
\label{sec:reconcile}
%-----------------------------------------------------------------------

Where both methods are applicable, they yield independent Monte Carlo
distributions. Reconciliation follows three rules based on the ratio
$\rho = m_{\mathrm{CO_2}}^{\mathrm{inv,p50}}\,/\,
m_{\mathrm{CO_2}}^{\mathrm{frp,p50}}$:

\textbf{Agreement} ($0.5 \le \rho \le 2.0$): both estimates are
published and the headline emission is the envelope distribution, formed
by pooling the Monte Carlo samples from both methods ($2N = 20{,}000$
draws) and computing p5/p50/p95 over the combined set. This approach
captures both within-method parametric uncertainty and between-method
structural uncertainty in a single distribution, without averaging away
genuine disagreement.

\textbf{Disagreement} ($\rho < 0.5$ or $\rho > 2.0$): both estimates
are retained but neither is published as a headline figure until editorial
review resolves the discrepancy. A written reviewer note must be entered
in the public changelog before publication.

\textbf{Third-party cross-check}: emission estimates from external sources
(e.g.\ CEOBS, CCI, or government statements) are recorded as
reference values only and do not enter headline arithmetic.

Events where $\rho$ falls within 10\% of either threshold
($\rho \in [0.50, 0.55] \cup [1.82, 2.00]$) are flagged for additional
editorial scrutiny even though they formally satisfy the agreement
criterion.

%-----------------------------------------------------------------------
\section{Verification and Confidence Assignment}
\label{sec:verify}
%-----------------------------------------------------------------------

Each candidate fire event is assessed against three independent evidence
streams:
\begin{itemize}
\item \textbf{Satellite persistence}: whether two or more independent
  satellite overpasses detect a hotspot at the same location.
\item \textbf{Optical confirmation}: whether a cloud-free Sentinel-2
  scene, acquired within $\pm$72 hours of the detected event, confirms
  active combustion.
\item \textbf{Conflict corroboration}: whether a documented armed event
  from a curated conflict database falls within a 24-hour, 2-kilometre
  spatio-temporal window of the detected fire.
\end{itemize}

Table~\ref{tab:confidence} defines the resulting confidence labels.

\begin{table}[h]
\centering
\caption{Confidence label decision table (v1.0.5).}
\label{tab:confidence}
\begin{tabular}{@{}lcccp{4.5cm}@{}}
\toprule
\textbf{Label} & \textbf{Persistent} &
\textbf{Optical} & \textbf{Conflict} & \textbf{Notes} \\
\midrule
\textsc{confirmed}  & yes & yes & yes & All three streams agree \\
\textsc{verified}   & yes & yes & no  & Satellite-only; no conflict record \\
\textsc{reported}   & yes & no  & any & Cloud cover or ambiguous scene \\
\textsc{suspected}  & no  & --- & any & Conflict record alone never auto-promotes \\
\textsc{suspected}  & no  & --- & no  & Single overpass only \\
\textsc{claimed}    & --- & --- & --- & News or official statement only; no satellite evidence \\
\bottomrule
\end{tabular}
\end{table}

The persistence threshold is two or more overpasses with elevated thermal
signal. Conflict-record corroboration alone is intentionally
non-promoting: a reported armed event near a facility without satellite
confirmation may represent a near-miss, a non-incendiary strike, or a
positional error in the conflict record. Such candidates are withheld from
publication pending review.

Only \textsc{confirmed} and \textsc{verified} events contribute to the
headline emission total. \textsc{reported} events appear on the public map
with a distinct marker. \textsc{suspected} and \textsc{claimed} events are
withheld pending review.

%-----------------------------------------------------------------------
\section{Uncertainty Quantification}
\label{sec:uq}
%-----------------------------------------------------------------------

Every emission estimate is the output of a Monte Carlo simulation with
$N = 10{,}000$ independent draws. Parameter prior distributions are
specified in Table~\ref{tab:priors}; they are not hard-coded in the
estimation logic and can be updated independently of it.

\begin{table}[h]
\centering
\caption{Monte Carlo parameter priors.}
\label{tab:priors}
\begin{tabular}{@{}llll@{}}
\toprule
\textbf{Parameter} & \textbf{Prior distribution} & \textbf{Units} &
\textbf{Source} \\
\midrule
Crude oil combustion factor
  & triangular(0.405, 0.425, 0.445)
  & tCO$_2$/barrel & EPA AP-42 \S1.3 \cite{ap42} \\
Refined product combustion factor
  & triangular(0.410, 0.430, 0.455)
  & tCO$_2$/barrel & IPCC 2006 GL Vol.~2 \cite{ipcc2006} \\
FRP-to-combustion rate, $\alpha$
  & $\mathcal{N}(0.368,\;0.05)$
  & kg/MJ & Wooster et al.~\cite{wooster2005} \\
Carbon recovery, $r$
  & triangular(0.92, 0.96, 0.98)
  & --- & Hobbs \& Radke~\cite{hobbs1992} \\
Burn duty cycle, $d$
  & triangular(0.40, 0.70, 0.95)
  & --- & Expert prior (Kuwait 1991; Abqaiq 2019) \\
Facility inventory at strike, $\phi$
  & uniform(0.30, 0.90)
  & --- & Judgment-based prior \\
FRP extrapolation factor, $k_{\mathrm{ext}}$
  & $\mathcal{N}(1.0,\;0.15)$
  & --- & MODIS/VIIRS cross-comparison \\
Fraction destroyed, $\psi$ (storage; default)
  & triangular(0.05, 0.15, 0.30)
  & --- & FRP/inventory ratio calibration (v1.0.5) \\
\bottomrule
\end{tabular}
\end{table}

\paragraph{Distribution semantics.}
Triangular distributions are parameterised as (low, mode, high). Where
asymmetric bounds are not explicitly specified, symmetric bounds around the
central value are assumed. Normal distributions use the stated standard
deviation; reported uncertainty bounds in the source literature are stored
for reviewer cross-check but are not used as truncation bounds.

\paragraph{Independence assumptions.}
Parameters are sampled independently for each event. Correlation of
parameters across facilities---for example, systematic satellite calibration
drift affecting all events simultaneously---is a planned extension for a
later version.

\paragraph{Aggregation.}
Facility-level and theatre-level aggregates are computed by element-wise
summation of per-event Monte Carlo sample arrays, then summarised to
p5/p50/p95. This approach correctly propagates uncertainty through the
summation rather than summing individual percentiles, which would
understate aggregate uncertainty by assuming all events simultaneously
attain their extreme values.

%-----------------------------------------------------------------------
\section{Methodology Version History}
\label{sec:versions}
%-----------------------------------------------------------------------

\begin{longtable}{@{}lllcp{6.5cm}@{}}
\toprule
\textbf{Version} & \textbf{Date} & \textbf{Type}
  & \textbf{Param.\ change} & \textbf{Summary} \\
\midrule
\endfirsthead
\toprule
\textbf{Version} & \textbf{Date} & \textbf{Type}
  & \textbf{Param.\ change} & \textbf{Summary} \\
\midrule
\endhead
\bottomrule
\endfoot
v1.0
  & 2026-03-01 & Initial & yes
  & Raw FRP method, inventory method, reconciliation, and confidence
    labels established. \\[4pt]
v1.0.1
  & 2026-05-24 & Calibration & yes
  & Baseline subtraction introduced: 75th-percentile rolling baseline
    with IQR-based robust standard deviation. Net FRP computed prior to
    Monte Carlo sampling. Conservative fallback baseline for events with
    insufficient pre-war observations. \\[4pt]
v1.0.2
  & 2026-05-24 & Data & no
  & Twelve months of pre-war archival satellite observations ingested.
    No equation or parameter change; baseline window now has sufficient
    historical coverage for all tracked facilities. \\[4pt]
v1.0.3
  & 2026-05-24 & Recompute & no
  & Full recomputation under v1.0.2 baselines. FRP/inventory agreement
    ratios of 2.7--7.4$\times$ observed for storage-type facilities,
    indicating systematic overestimation of fraction destroyed. No
    parameter change. \\[4pt]
v1.0.4
  & 2026-05-24 & Recompute & no
  & Confirmed systematic overestimation with additional data. No parameter
    change. \\[4pt]
v1.0.5
  & 2026-05-24 & Calibration & yes
  & Fraction-destroyed defaults for storage-type facilities recalibrated
    to triangular(0.05, 0.15, 0.30). Damage assessment records updated
    to reflect recalibrated defaults. \\
\bottomrule
\end{longtable}

\noindent This document supersedes methodology v1.0 (2026-05-23). All
equation numbers are unchanged relative to v1.0; material changes are in
Section~\ref{sec:frp} (baseline subtraction made explicit),
Section~\ref{sec:inventory} (facility-type restriction and
fraction-destroyed recalibration), Section~\ref{sec:reconcile} (envelope
distribution defined), and this section (version history added).

%-----------------------------------------------------------------------
\section{Worked Example: Ahvaz/Karoon, 2026-02-28}
\label{sec:example}
%-----------------------------------------------------------------------

\paragraph{Context.}
As of 2026-05-24, the system has processed 47 fire events across 7 Iranian
facilities under methodology v1.0.5. Of these, 27 events have FRP-based
emission estimates; a subset also carry inventory-based estimates. Two
events are fully reconciled ($\rho \in [0.5, 2.0]$); none are flagged as
requiring review. The cumulative headline is
${\sim}75{,}619$~tCO$_2$e (p50) across all tracked facilities.

\paragraph{Event.}
The Ahvaz/Karoon production area event of 2026-02-28 is the largest single
event in the current dataset. Its inputs are:
\begin{itemize}
\item Peak FRP: 114~MW.
\item Duration: 6.5~days.
\item Raw FRP integral: $I_{\mathrm{raw}} = 2.893 \times 10^{7}$~MJ
  ($\approx 28{,}930$~GJ), from VIIRS and MODIS overpasses.
\item Pre-war baseline: $P_{75} = 13.59$~MW (from 12 months of archival
  observations at this facility).
\item Destroyed-fraction assessment: triangular(0.05, 0.15, 0.30)
  (v1.0.5 default; no reviewer-assessed $\hat\psi$ available).
\item Confidence label: \textsc{reported} (satellite-persistent but no
  cloud-free optical confirmation obtained for this event window).
\end{itemize}

\paragraph{FRP-based point estimate.}
The following uses central prior values
($k_{\mathrm{ext}} = 1.0$, $d = 0.70$, $\alpha = 0.368$~kg/MJ) to
illustrate the equation chain before Monte Carlo sampling.

\noindent\emph{Stage 2 --- Baseline subtraction}
(Eq.~\eqref{eq:net-frp}):
\begin{align*}
I_{\mathrm{net}}
  &= I_{\mathrm{raw}}
     - \mathrm{FRP}_{\mathrm{baseline}} \cdot \Delta t \\
  &= 2.893 \times 10^{7}
     - 13.59\;\mathrm{MW}
       \times 6.5\;\mathrm{d}
       \times 86{,}400\;\mathrm{s\,d^{-1}} \\
  &= 2.893 \times 10^{7} - 7.634 \times 10^{6} \\
  &= 2.130 \times 10^{7}\;\text{MJ}.
\end{align*}

\noindent\emph{Stage 3 --- Overpass-gap correction}
(Eq.~\eqref{eq:fre}):
\begin{align*}
I_{\mathrm{FRE}}
  &= k_{\mathrm{ext}} \cdot d \cdot I_{\mathrm{net}}
   = 1.0 \times 0.70 \times 2.130 \times 10^{7}
   = 1.491 \times 10^{7}\;\text{MJ}.
\end{align*}

\noindent\emph{Stage 4 --- Conversion to $\mathrm{CO_2}$ mass}
(Eqs.~\eqref{eq:fuel} and~\eqref{eq:co2-frp}):
\begin{align*}
m_{\mathrm{fuel}}
  &= 0.368 \times 1.491 \times 10^{7}
   = 5{,}487\;\text{t fuel}, \\
m_{\mathrm{CO_2}}^{\mathrm{frp}}
  &= 5{,}487 \times 3.15
   = 17{,}284\;\text{tCO}_2\text{e}.
\end{align*}

\paragraph{Monte Carlo result.}
The point estimate (17,284~tCO$_2$e) falls between the fifth percentile
and median of the full Monte Carlo distribution, as expected: the right
skew in $d$ and $k_{\mathrm{ext}}$ pulls the median above the product of
central values. The realised distribution from the live system is:

\begin{center}
\begin{tabular}{@{}ccc@{}}
\toprule
p5 & p50 & p95 \\
\midrule
10,679~tCO$_2$e & 22,704~tCO$_2$e & 55,788~tCO$_2$e \\
\bottomrule
\end{tabular}
\end{center}

\noindent This event is classified as \textsc{reported} rather than
\textsc{verified} or \textsc{confirmed} because no cloud-free Sentinel-2
scene was obtained for the 28~February event window. Its contribution to
the headline dashboard total is included, but its higher confidence
threshold will be revisited if optical imagery becomes available.

%-----------------------------------------------------------------------
\begin{thebibliography}{14}
%-----------------------------------------------------------------------

\bibitem{ap42}
  U.S.\ Environmental Protection Agency.
  \emph{AP-42: Compilation of Air Emissions Factors},
  Section~1.3 (fuel oil combustion). EPA.

\bibitem{ipcc2006}
  IPCC (2006).
  \emph{2006 IPCC Guidelines for National Greenhouse Gas Inventories},
  Volume~2, Chapter~1, Table~1.3. IPCC.

\bibitem{wooster2005}
  Wooster, M.J., Roberts, G., Perry, G.L.W., Kaufman, Y.J. (2005).
  Retrieval of biomass combustion rates and totals from fire radiative
  power observations.
  \emph{Journal of Geophysical Research}, 110, D24311.

\bibitem{hobbs1992}
  Hobbs, P.V., Radke, L.F. (1992).
  Airborne studies of the smoke from the Kuwait oil fires.
  \emph{Science}, 256(5059), 987--991.

\bibitem{cci2026}
  Climate and Community Institute (2026).
  \emph{Iran 14-day conflict emissions brief}. CCI.

\bibitem{ceobs}
  Conflict and Environment Observatory (2026).
  Incident records, Iran conflict.
  \url{https://ceobs.org}.

\bibitem{elvidge2016}
  Elvidge, C.D., Zhizhin, M., Hsu, F.-C., Baugh, K. (2016).
  Methods for global survey of natural gas flaring from Visible Infrared
  Imaging Radiometer Suite data.
  \emph{Energies}, 9(1), 14.

\bibitem{freeborn2014}
  Freeborn, P.H., Wooster, M.J., Roberts, G. (2014).
  Addressing the spatiotemporal sampling design of MODIS to provide
  estimates of the fire radiative energy emitted from Africa.
  \emph{Remote Sensing of Environment}, 140, 697--734.

\bibitem{neimark2026}
  Neimark, B., Otu-Larbi, F., Bigger, P. (2026).
  Carbon emissions from the Gaza war.
  \emph{One Earth}.

\bibitem{deklerk2024}
  De Klerk, L., Shlapak, M., Krakovska, S.\ et al.\ (2024).
  \emph{Guidance on the Assessment of Conflict-Related GHG Emissions,
  v1.0}. Initiative on GHG Accounting of War / Ecoaction.

\bibitem{bigger2026}
  Bigger, P., Neimark, B., Otu-Larbi, F. (2026).
  Two weeks of war in Iran unleashed more carbon pollution than Iceland.
  Climate and Community Institute Research Snapshot.

\bibitem{cottrell2022}
  Cottrell, L. (2022).
  A framework for military greenhouse gas emissions reporting.
  Conflict and Environment Observatory (CEOBS).

\bibitem{vandenhoek2024}
  Van Den Hoek, J., Scher, C. (2024).
  Sentinel-1 SAR coherent change detection for conflict damage assessment.
  Oregon State University Conflict Ecology Lab.

\bibitem{crawford2019}
  Crawford, N. (2019).
  \emph{Pentagon Fuel Use, Climate Change, and the Costs of War}.
  Costs of War Project, Brown University.

\end{thebibliography}

\end{document}